\begin{center}
\LARGE
Building abelian groups and commutative rings into
                 first-order theorem proving
\end{center}

\begin{center}
\Large
J\"urgen Stuber
\end{center}

\noindent
Date:  July 17th (Wednesday)\\
Time:  14:00-14:25\\

\begin{center}
\large
Abstract
\end{center}

For many theories from algebra there exist convergent term rewriting
systems, possibly modulo some equational theory. We present a
technique to build such a theory into ordering-restricted first-order
clausal theorem proving. In this way one obtains a refutationally
complete inference system which creates a smaller search space than by
directly applying a more general prover to the axioms of the theory.

Specifically, on the ground level equational literals are simplified,
until they are in a {\it sufficiently} {\it simplified} form. For each of these
forms one derives a {\it symmetrization}, that is a set of rules such that
critical pairs with the theory converge. One then considers overlaps
between symmetrizations.  Since their structure is know, one can
exploit this by showing that only a few critical pairs need to be
tested to obtain convergence.

To lift this to the non-ground level, each step of the simplification
becomes an inference. These are either carefully controlled
superpositions with theory axioms or derived inferences like
cancellation. Then, for inferences between sufficiently simplified
non-theory clauses, one distinguished between overlaps of two extended
rules and overlaps where at least one rule is not extended. The former
correspond to the critical pairs obtained above; the latter are
similar to superposition inferences, except that they take the
extended rules in the symmetrization into account.

The technique is quite general, but it seems most useful for theories
which have abelian groups as a subtheory. In that case a single term
of a sum can be isolated on the left-hand side. The inference systems
obtained are extensions of the corresponding Gr�bner base
algorithms.

